%% ****** Start of file apstemplate.tex ****** %
%%
%%
%%   This file is part of the APS files in the REVTeX 4.2 distribution.
%%   Version 4.2a of REVTeX, January, 2015
%%
%%
%%   Copyright (c) 2015 The American Physical Society.
%%
%%   See the REVTeX 4 README file for restrictions and more information.
%%
%
% This is a template for producing manuscripts for use with REVTEX 4.2
% Copy this file to another name and then work on that file.
% That way, you always have this original template file to use.
%
% Group addresses by affiliation; use superscriptaddress for long
% author lists, or if there are many overlapping affiliations.
% For Phys. Rev. appearance, change preprint to twocolumn.
% Choose pra, prb, prc, prd, pre, prl, prstab, prstper, or rmp for journal
%  Add 'draft' option to mark overfull boxes with black boxes
%  Add 'showkeys' option to make keywords appear
%\documentclass[aps,prl,preprint,groupedaddress]{revtex4-2}
\documentclass[aps,prl,twocolumn,groupedaddress]{revtex4-2}
%\documentclass[aps,prl,preprint,superscriptaddress]{revtex4-2}
%\documentclass[aps,prl,reprint,groupedaddress]{revtex4-2}

% You should use BibTeX and apsrev.bst for references
% Choosing a journal automatically selects the correct APS
% BibTeX style file (bst file), so only uncomment the line
% below if necessary.
%\bibliographystyle{apsrev4-2}

\usepackage{graphicx}
\usepackage{epstopdf}
\usepackage{caption}
%\usepackage{amsmath}% http://ctan.org/pkg/amsmath
%\usepackage{amsthm}
%\usepackage{amsfonts}
%\usepackage{subfigure}
%\usepackage{hhline}
%\usepackage[miktex]{gnuplottex}
%\usepackage{xcolor}
\usepackage{amssymb}
\usepackage{amsmath}
\usepackage{color}
\usepackage{hyperref}
%\usepackage[percent]{overpic}
\usepackage{tikz}
\usepackage{mathrsfs}
\usepackage{wasysym}
\usepackage{tikz-cd}
%\usepackage{linearb}
%\usepackage{wsuipa}
\usepackage{calc}
\usepackage{tipa}

%\usepackage{stix} %\fisheye
\usepackage{stackengine,scalerel}

% so sections, subsections, etc. become numerated.
\setcounter{secnumdepth}{3}

\newcommand{\FF}{\mathbb{F}}
\newcommand{\NN}{\mathbb{N}}
\newcommand{\ZZ}{\mathbb{Z}}
\newcommand{\QQ}{\mathbb{Q}}
\newcommand{\RR}{\mathbb{R}}
\newcommand{\CC}{\mathbb{C}}
\newcommand{\II}{\mathbb{I}}
\newcommand{\GG}{\mathbb{G}}

\DeclareMathOperator*{\argmax}{arg\,max}
\DeclareMathOperator*{\argmin}{arg\,min}
\newcommand{\avrg}[1]{\left\langle #1 \right\rangle}
\newcommand{\nelta}{\bar{\delta}}
\newcommand{\bra}[1]{\left\langle #1\right|}
\newcommand{\ket}[1]{\left| #1 \right\rangle}
\newcommand{\sbra}[1]{\langle #1|}
\newcommand{\sket}[1]{| #1 \rangle}
\newcommand{\bek}[3]{\left\langle #1 \right| #2 \left| #3 \right\rangle}
\newcommand{\sbek}[3]{\langle #1 | #2 | #3 \rangle}
\newcommand{\braket}[2]{\left\langle #1 \middle| #2 \right\rangle}
\newcommand{\ketbra}[2]{\left| #1 \middle\rangle \middle\langle #2  \right|}
\newcommand{\sbraket}[2]{\langle #1 | #2 \rangle}
\newcommand{\sketbra}[2]{| #1 \rangle  \langle #2 |}
\newcommand{\norm}[1]{\left\lVert#1\right\rVert}
\newcommand{\snorm}[1]{\lVert#1\rVert}
\newcommand{\bvec}[1]{\boldsymbol{\mathsf{#1}}}
\newcommand{\bcov}[1]{\boldsymbol{#1}}
\newcommand{\bdua}[1]{\boldsymbol{\check{#1}}}
%\newcommand{\bdov}[1]{\boldsymbol{\breve{#1}}}
\newcommand{\bdov}[1]{\breve{#1}}
%\newcommand{\bten}[1]{\boldsymbol{\mathfrak{#1}}}
\newcommand{\bten}[1]{\boldsymbol{\mathfrak{#1}}}
\newcommand{\forany}{\tilde{\forall}}
\newcommand{\qed}{$\overset{\circ}{.}\;$}
\newcommand{\bigo}{\mathcal{O}}
\newcommand{\spn}{\mathrm{spn}\,}
\newcommand{\krn}{\mathrm{krn}\,}
\newcommand{\biy}{\leftrightarrow}
\newcommand{\ugh}{\mathbin{\textlinb{\BPwheel}}}

\newcommand\bigeye{\ensurestackMath{\stackinset{c}{}{c}{-.3pt}%
  {\bullet}{\scriptstyle\bigcirc}}}
\newcommand\eye{\scalerel*{\bigeye}{x}}
%\newcommand*{\fisheye}{%
%    \mathbin{%
%        \ooalign{$\circledcirc$\cr\hidewidth$\bullet$\hidewidth}%
%    }%
%}

% double angle-bracket
% https://tex.stackexchange.com/questions/79657/how-to-get-double-angle-bracket-without-using-mnsymbol-package
\makeatletter
\newsavebox{\@brx}
\newcommand{\llangle}[1][]{\savebox{\@brx}{\(\m@th{#1\langle}\)}%
  \mathopen{\copy\@brx\mkern2mu\kern-0.9\wd\@brx\usebox{\@brx}}}
\newcommand{\rrangle}[1][]{\savebox{\@brx}{\(\m@th{#1\rangle}\)}%
  \mathclose{\copy\@brx\mkern2mu\kern-0.9\wd\@brx\usebox{\@brx}}}
\makeatother

% https://tex.stackexchange.com/questions/297362/bar-through-letter-in-mathmode
%\usepackage{calc}
\newsavebox\CBox
\newcommand\hcancel[2][0.5pt]{%
  \ifmmode\sbox\CBox{$#2$}\else\sbox\CBox{#2}\fi%
  \makebox[0pt][l]{\usebox\CBox}%  
  \rule[0.77\ht\CBox-#1/2]{\wd\CBox}{#1}}
\newcommand{\cb}{\hcancel{b}}
\newcommand{\cd}{\hcancel{d}}
\newcommand{\ct}{\hcancel{\partial}}

\begin{document}

% Use the \preprint command to place your local institutional report
% number in the upper righthand corner of the title page in preprint mode.
% Multiple \preprint commands are allowed.
% Use the 'preprintnumbers' class option to override journal defaults
% to display numbers if necessary
%\preprint{}

%Title of paper
\title{
M\'etodos Matem\'aticos de la F\'isica II
}

% repeat the \author .. \affiliation  etc. as needed
% \email, \thanks, \homepage, \altaffiliation all apply to the current
% author. Explanatory text should go in the []'s, actual e-mail
% address or url should go in the {}'s for \email and \homepage.
% Please use the appropriate macro foreach each type of information

% \affiliation command applies to all authors since the last
% \affiliation command. The \affiliation command should follow the
% other information
% \affiliation can be followed by \email, \homepage, \thanks as well.
\author{Juan I. Perotti}
\email[]{juan.perotti@unc.edu.ar}
%\homepage[]{Your web page}
%\thanks{}
%\altaffiliation{}
%\affiliation{}
\affiliation{Instituto de F\'isica Enrique Gaviola (IFEG-CONICET), Ciudad Universitaria, 5000 C\'ordoba, Argentina}
\affiliation{Facultad de Matem\'atica, Astronom\'ia, F\'isica y Computaci\'on, Universidad Nacional de C\'ordoba, Ciudad Universitaria, 5000 C\'ordoba, Argentina}

% \author{Colaborador}
% \affiliation{Instituto de F\'isica Enrique Gaviola (IFEG-CONICET), Ciudad Universitaria, 5000 C\'ordoba, Argentina}
% \affiliation{Facultad de Matem\'atica, Astronom\'ia, F\'isica y Computaci\'on, Universidad Nacional de C\'ordoba, Ciudad Universitaria, 5000 C\'ordoba, Argentina}
% \email[E-mail: ]{xxx.yyy@unc.edu.ar }

%Collaboration name if desired (requires use of superscriptaddress
%option in \documentclass). \noaffiliation is required (may also be
%used with the \author command).
%\collaboration can be followed by \email, \homepage, \thanks as well.
%\collaboration{Claudio J. Tessone}
%\noaffiliation

\date{\today}

\begin{abstract}
Bla bla bla... resumen...
\end{abstract}

% insert suggested keywords - APS authors don't need to do this
%\keywords{}

%\maketitle must follow title, authors, abstract, and keywords
\maketitle
\onecolumngrid

\section{Campos Tensoriales y Cambios de Coordenadas}

Esta secci\'on son notas de lectura de las notas escritas por Omar Ortiz que pueden encontrarse en el archivo \verb+metodos-matematicos-de-la-fisica-II-2026-priv/aula-virtual-2025/notas-de-clase/Campos_tensoriales-Cambios_de_coordenadas.pdf+.
En particular, se resuelven algunos de los ejercicios planteados en las notas de Omar.

\subsection{Ejercicio 1}

% https://chatgpt.com/c/69d9505f-7538-83e9-907a-7446aafe082a

\subsubsection{Forma 1 de pensarlo}

Considere un paraboloide hiperb\'olico embebido en el espacio Eucl\'ideo tridimensional dado, en coordenadas cartesianas, por $z=x^2-y^2$.
Escriba la m\'etrica del espacio Eucl\'ideo en coordenadas $\{s,t,u\}$ dadas por 
$$
s=x,\,\, t=y,\,\, u=z+x^2-y^2.
$$
Muestre ahora que la m\'etrica $h_{ij}$ inducida sobre el paraboloide hiperb\'olico es
$$
h_{11}=1+4s^2,\,\,
h_{12}=h_{21} = -4st,\,\,
h_{22}=1+4t^2.
$$

Si invertimos el cambio de coordenadas
$$
s(x,y,z)=x
$$
$$
t(x,y,z)=y
$$
$$
u(x,y,z)=z+x^2-y^2
$$
obtenemos
$$
x(s,t,u)=s
$$
$$
y(s,t,u)=t
$$
$$
z(s,t,u)=u-s^2+t^2
.
$$
Luego, el cuadrado de una distancia infinitesimal viene dado por
\begin{eqnarray}
d\ell^2
=
dx^2
+
dy^2
+
dz^2
&=&
ds^2
+
dt^2
+
4s^2\,ds^2
+
4t^2\,dt^2
+
du^2
-
4st\,ds\,dt
-
4st\,dt\,ds
-
2s\,ds\,du
+
2t\,dt\,du
+
\notag
\\
&=&
(1+4s^2)ds^2
+
(1+4t^2)dt^2
-
4st\,ds\,dt
-
4st\,dt\,ds
+
\bigo(du)
\notag
\\
&=&
g_{ss}ds^2
+
g_{tt}dt^2
+
g_{st}\,ds\,dt
+
g_{ts}\,dt\,ds
+
\bigo(du)
\notag
\end{eqnarray}
puesto que
$$
dz
=
\frac{\partial z}{\partial s}\bigg|_{t,u}
ds
+
\frac{\partial z}{\partial t}\bigg|_{s,u}
dt
+
\frac{\partial z}{\partial u}\bigg|_{s,t}
du
=
-2s\,ds
+2t\,dt
+du
$$
y, por ende
\begin{eqnarray}
dz^2
&=&
\big(
-2s\,ds
+2t\,dt
+du
\big)^2
\notag
\\
&=&
4s^2\,ds^2
+
4t^2\,dt^2
+
du^2
-
4st\,ds\,dt
-
4st\,dt\,ds
-
2s\,ds\,du
+
2t\,dt\,du
\notag
\\
&=&
4s^2\,ds^2
+
4t^2\,dt^2
-
4st\,ds\,dt
-
4st\,dt\,ds
+
\bigo(du)
.
\notag
\end{eqnarray}
Todo lo anterior determina unívocamente las componentes de la m\'etrica $g$ en las coordenadas $s$, $t$ y $u$.
En un desplazamiento infinitesimal por una superficie determinada por la restricci\'on $u=\mathrm{constante}$, se tiene que $du=0$.
M\'as precisamente, se tiene que $du(v)=0$ para todo vector tangente $v$ perteneciente al espacio tangente a la subvariedad definida por la superficie a $u=\mathrm{constante}$ en el punto arbitrario de la misma que está en consideración.
Por ende, se obtiene la m\'etrica reducida de coordenadas
$$
h_{ss} = g_{ss} = 1+4s^2
$$
$$
h_{tt} = g_{tt} = 1+4t^2
$$
y
$$
h_{st} = g_{st} = -4st
.
$$
Notar que hay un pequeño error en las notas, en donde incorrectamente dice que $h_{st}=-2st$.

\subsubsection{Forma 2 de pensarlo}

Si una hipersuperficie en $\RR^n$ está caracterizada por coordenadas $k^1,...,k^m$ con $m\leq n$, la métrica inducida viene dada por
$$
h_{ab}
=
g_{ij}
\frac{\partial q^i}{\partial k^a}
\frac{\partial q^j}{\partial k^b}
$$
donde $q^1,...,q^n$ son las coordenadas en $\RR^n$ y 
$$
g_{ij} 
= 
\bigg\langle 
\frac{\partial}{\partial q_i} 
,
\frac{\partial}{\partial q_j} 
\bigg\rangle
$$ 
es la métrica asociada correspondiente.

En el presente problema, tenemos que
$$
q^1 = x,\, 
q^2 = y,\, 
q^3 = z
$$
$$
g_{ij} = \delta_{ij}
$$
$$
k^1 = s,\,
k^2 = t,\,
k^3 = u
.
$$
Además, la superficie viene determinada por las condiciones:
\begin{itemize}
\item  $u=\mathrm{constante}=0$, y
\item $z = u - s^2 + t^2.$
\end{itemize}
En otras palabras, la superficie viene determinada por el mapa $\RR^2\ni (s,t)\to (x,y,z)\in \RR^3$ definido por las expresiones:
$$
x(s,t,u=u_0=0) = s
$$
$$
y(s,t,u=u_0=0) = t
$$
$$
z(s,t,u=u_0=0) = u_0 - x^2(s,t,u=u_0=0) + y^2(s,t,u=u_0=0) = 0 - s^2 + t^2
.
$$
De esta manera, la métrica inducida tiene componentes
\begin{eqnarray}
h_{ss}
&=&
g_{xx}
\frac{\partial x}{\partial s}
\frac{\partial x}{\partial s}
+
g_{xy}
\frac{\partial x}{\partial s}
\frac{\partial y}{\partial s}
+
g_{xz}
\frac{\partial x}{\partial s}
\frac{\partial z}{\partial s}
\notag
\\
&&
+
g_{yx}
\frac{\partial y}{\partial s}
\frac{\partial x}{\partial s}
+
g_{yy}
\frac{\partial y}{\partial s}
\frac{\partial y}{\partial s}
+
g_{yz}
\frac{\partial y}{\partial s}
\frac{\partial z}{\partial s}
\notag
\\
&&
+
g_{zx}
\frac{\partial z}{\partial s}
\frac{\partial x}{\partial s}
+
g_{zy}
\frac{\partial z}{\partial s}
\frac{\partial y}{\partial s}
+
g_{zz}
\frac{\partial z}{\partial s}
\frac{\partial z}{\partial s}
\notag
\\
&=&
1 + 0 + 0 \notag\\
&&
+ 0 + 0 + 0 \notag\\
&&
+ 0 + 0 + (2s)^2
\notag
\\
&=&
1 + 4s^2.
\notag
\end{eqnarray}
\begin{eqnarray}
h_{st}
&=&
g_{xx}
\frac{\partial x}{\partial s}
\frac{\partial x}{\partial t}
+
g_{xy}
\frac{\partial x}{\partial s}
\frac{\partial y}{\partial t}
+
g_{xz}
\frac{\partial x}{\partial s}
\frac{\partial z}{\partial t}
\notag
\\
&&
+
g_{yx}
\frac{\partial y}{\partial s}
\frac{\partial x}{\partial t}
+
g_{yy}
\frac{\partial y}{\partial s}
\frac{\partial y}{\partial t}
+
g_{yz}
\frac{\partial y}{\partial s}
\frac{\partial z}{\partial t}
\notag
\\
&&
+
g_{zx}
\frac{\partial z}{\partial s}
\frac{\partial x}{\partial t}
+
g_{zy}
\frac{\partial z}{\partial s}
\frac{\partial y}{\partial t}
+
g_{zz}
\frac{\partial z}{\partial s}
\frac{\partial z}{\partial t}
\notag
\\
&=&
0 + 0 + 0 \notag\\
&&
+ 0 + 0 + 0 \notag\\
&&
+ 0 + 0 - 2s \, 2t
\notag
\\
&=&
-4st.
\notag
\end{eqnarray}
\begin{eqnarray}
h_{tt}
&=&
g_{xx}
\frac{\partial x}{\partial t}
\frac{\partial x}{\partial t}
+
g_{xy}
\frac{\partial x}{\partial t}
\frac{\partial y}{\partial t}
+
g_{xz}
\frac{\partial x}{\partial t}
\frac{\partial z}{\partial t}
\notag
\\
&&
+
g_{yx}
\frac{\partial y}{\partial t}
\frac{\partial x}{\partial t}
+
g_{yy}
\frac{\partial y}{\partial t}
\frac{\partial y}{\partial t}
+
g_{yz}
\frac{\partial y}{\partial t}
\frac{\partial z}{\partial t}
\notag
\\
&&
+
g_{zx}
\frac{\partial z}{\partial t}
\frac{\partial x}{\partial t}
+
g_{zy}
\frac{\partial z}{\partial t}
\frac{\partial y}{\partial t}
+
g_{zz}
\frac{\partial z}{\partial t}
\frac{\partial z}{\partial t}
\notag
\\
&=&
1 + 0 + 0 \notag\\
&&
+ 0 + 0 + 0 \notag\\
&&
+ 0 + 0 + (-2t)^2
\notag
\\
&=&
1 + 4t^2.
\notag
\end{eqnarray}



%See fig.~\ref{fig1}.
%
%\begin{figure*}
%\includegraphics*[scale=.4]{assets/fig1.png}
%\includegraphics*[scale=.4]{assets/fig1.png}
%\put(-493,140){\bf a)}
%\put(-250,140){\bf b)}
%\\
%\includegraphics*[scale=.4]{assets/fig1.png}
%\includegraphics*[scale=.4]{assets/fig1.png}
%\put(-493,140){\bf c)}
%\put(-250,140){\bf d)}
%\vspace{-0.25cm}
%\caption{
%\label{fig1}
%Bla bla.
%{\bf a)} 
%Bla bla.
%{\bf b)} 
%Bla bla.
%{\bf c)} 
%Bla bla.
%{\bf d)} 
%Bla bla.
%}
%\end{figure*}

% If in two-column mode, this environment will change to single-column
% format so that long equations can be displayed. Use
% sparingly.
%\begin{widetext}
% put long equation here
%\end{widetext}

% figures should be put into the text as floats.
% Use the graphics or graphicx packages (distributed with LaTeX2e)
% and the \includegraphics macro defined in those packages.
% See the LaTeX Graphics Companion by Michel Goosens, Sebastian Rahtz,
% and Frank Mittelbach for instance.
%
% Here is an example of the general form of a figure:
% Fill in the caption in the braces of the \caption{} command. Put the label
% that you will use with \ref{} command in the braces of the \label{} command.
% Use the figure* environment if the figure should span across the
% entire page. There is no need to do explicit centering.

% \begin{figure}
% \includegraphics{}%
% \caption{\label{}}
% \end{figure}

% Surround figure environment with turnpage environment for landscape
% figure
% \begin{turnpage}
% \begin{figure}
% \includegraphics{}%
% \caption{\label{}}
% \end{figure}
% \end{turnpage}

% tables should appear as floats within the text
%
% Here is an example of the general form of a table:
% Fill in the caption in the braces of the \caption{} command. Put the label
% that you will use with \ref{} command in the braces of the \label{} command.
% Insert the column specifiers (l, r, c, d, etc.) in the empty braces of the
% \begin{tabular}{} command.
% The ruledtabular enviroment adds doubled rules to table and sets a
% reasonable default table settings.
% Use the table* environment to get a full-width table in two-column
% Add \usepackage{longtable} and the longtable (or longtable*}
% environment for nicely formatted long tables. Or use the the [H]
% placement option to break a long table (with less control than 
% in longtable).
% \begin{table}%[H] add [H] placement to break table across pages
% \caption{\label{}}
% \begin{ruledtabular}
% \begin{tabular}{}
% Lines of table here ending with \\
% \end{tabular}
% \end{ruledtabular}
% \end{table}

% Surround table environment with turnpage environment for landscape
% table
% \begin{turnpage}
% \begin{table}
% \caption{\label{}}
% \begin{ruledtabular}
% \begin{tabular}{}
% \end{tabular}
% \end{ruledtabular}
% \end{table}
% \end{turnpage}

%\begin{acknowledgments}
%JIP acknowledges financial support from grants CONICET (PIP 112 20150 10028), FonCyT (PICT-2017-0973), SeCyT–UNC (Argentina) and MinCyT C\'ordoba (PID PGC 2018).
%The authors thank CCAD -- Universidad Nacional de C\'ordoba, \url{http://ccad.unc.edu.ar/}, which is part of SNCAD -- MinCyT, Argentina, for the provided computational resources.
%\end{acknowledgments}
%% Create the reference section using BibTeX:
%\bibliography{ref}

% Specify following sections are appendices. Use \appendix* if there
% only one appendix.

%\onecolumngrid

\appendix

\section{Definiciones alternativas de espacio tangente}

Basado en ChatGPT-2026-04-26:
\url{https://chatgpt.com/share/69e8bec0-b0a4-8325-a5d5-6a045a4fdb2c}.

For a smooth manifold $M$ and point $p\in M$, the tangent space $T_pM$ can be defined in several equivalent ways.

\subsection{Tangent vectors as velocities of curves (Geometric definition)}

A tangent vector at (p) is an equivalence class of smooth curves

$$
\gamma:(-\epsilon,\epsilon)\to M
$$

with

$$
\gamma(0)=p.
$$

Two curves are equivalent if they have the same first-order behavior at (p), meaning in any (equivalently every) chart

$$
\frac{d(x^i\circ \gamma_1)}{dt}(0)
=
\frac{d(x^i\circ \gamma_2)}{dt}(0).
$$

Then

$$
v=\gamma'(0)
$$

is a tangent vector.

So

$$
T_pM
=
{\text{velocities of curves through }p}.
$$

\subsubsection{Intuition:}

Tangent vectors are infinitesimal motions.

\subsection{Tangent vectors as derivations (Operator definition)}

A tangent vector is a linear map

$$
v:C^\infty(M)\to \mathbb R
$$

such that

$$
v(fg)
=
f(p)v(g)+g(p)v(f).
$$

(Leibniz rule.)

These are derivations at $p$.

Then

$$
T_pM
=
{\text{all derivations at }p}.
$$

Coordinate basis:

$$
\frac{\partial}{\partial x^i}.
$$

General vector:

$$
v=v^i \frac{\partial}{\partial x^i}.
$$

\subsubsection{Intuition:}

Tangent vectors are directional derivatives.

\subsection{Tangent vectors as coordinate tuples (Local coordinate definition)}

Choose coordinates near $p$:

$$
(x^1,\dots,x^n).
$$

Then a tangent vector is just

$$
(v^1,\dots,v^n)\in \mathbb R^n.
$$

with basis

$$
\frac{\partial}{\partial x^i}.
$$

Thus

$$
T_pM\simeq \mathbb R^n
$$

locally.

This is often how physicists first meet tangent spaces.

\subsection{Tangent vectors as equivalence classes of germs (More algebraic)}

Take germs of functions vanishing at $p$, let

$$
\mathfrak m_p
$$

be the maximal ideal.

Then

$$
T_pM
\cong
(\mathfrak m_p/\mathfrak m_p^2)^*.
$$

This says tangent vectors are dual to first-order infinitesimals.

Very elegant algebraic definition.

\subsection{Tangent vectors as pushforwards of coordinate vectors (Embedded viewpoint)}

If

$$
r(u^1,\dots,u^n)
$$

parametrizes a manifold, then tangent vectors are spanned by

$$
\frac{\partial r}{\partial u^i}.
$$

Thus

$$
T_pM
=
\text{span}
\{
\partial_i r
\}
.
$$

This is what Briozzo used in Problem 4.

For surfaces in $\mathbb R^3$, this is the most concrete picture.

\subsection{Tangent space as best linear approximation (Differential viewpoint)}

A manifold near $p$ looks infinitesimally like a vector space:

$$
T_pM
$$

is the linear space approximating $M$ near $p$.

This is the ``linearization'' viewpoint.

\subsection{If $M\subset \mathbb R^n$ is given implicitly (Constraint definition)}

Suppose

$$
M={x:F(x)=0}
$$

with

$$
F:\mathbb R^n\to \mathbb R^k.
$$

Then

$$
T_pM
=
\ker dF_p.
$$

Equivalent:

vectors satisfying linearized constraints.

Example:
Sphere

$$
x^2+y^2+z^2=1
$$

gives

$$
2xv_x+2yv_y+2zv_z=0.
$$

So tangent vectors are orthogonal to radius.

Very important in physics.

\subsection{As dual of cotangent space (Covector duality)}

Define cotangent space first:

$$
T_p^*M.
$$

Then

$$
T_pM=(T_p^*M)^*.
$$

Especially natural in modern geometry.

\subsection{All equivalent}

These all define the same object.
There are natural isomorphisms:
Curve velocity
$
\longleftrightarrow
$
Derivation
$
\longleftrightarrow
$
Coordinate vector
$
\longleftrightarrow
$
Embedded tangent arrow
$
\longleftrightarrow
$
Kernel of constraint map
$
\longleftrightarrow
$
Dual of infinitesimals.
All these are different representations of the same tangent vector.

Summary table:

%| Viewpoint                | Tangent vector is...             |
%| ------------------------ | -------------------------------- |
%| Curves                   | Velocity class (\gamma'(0))      |
%| Derivations              | Directional derivative           |
%| Coordinates              | Tuple ((v^1,\dots,v^n))          |
%| Parametrized submanifold | Span of (\partial_i r)           |
%| Implicit manifold        | Kernel of (dF)                   |
%| Algebraic                | ((\mathfrak m/\mathfrak m^2)^*)  |
%| Linearization            | Best infinitesimal approximation |

\begin{table}[h]
\centering
\begin{tabular}{|l|l|}
\hline
\textbf{Viewpoint} & \textbf{Tangent vector is...} \\
\hline
Curves & Velocity class $\gamma'(0)$ \\
Derivations & Directional derivative \\
Coordinates & Tuple $(v^1,\dots,v^n)$ \\
Parametrized submanifold & Span of $\partial_i r$ \\
Implicit manifold & Kernel of $dF$ \\
Algebraic & $(\mathfrak m/\mathfrak m^2)^*$ \\
Linearization & Best infinitesimal approximation \\
\hline
\end{tabular}
\caption{Equivalent viewpoints on tangent vectors.}
\end{table}

My recommendation (for physics). Keep these three always linked:

$$
\boxed{
\text{velocity}
\leftrightarrow
\text{derivation}
\leftrightarrow
\text{embedded arrow}
}
$$

\subsection{Proof of equivalences}

The fundamental theorem is usually proving equivalence between the three core definitions:
\begin{itemize}
\item Curves.
\item Derivations.
\item Coordinate vectors ($\mathbb R^n$).
\end{itemize}
Once those are equivalent, the embedded and implicit definitions follow fairly naturally.

\vspace{0.5cm}
{\bf Theorem:} For an $n$-dimensional smooth manifold $M$ and $p\in M$,
$$
T_pM
$$
defined via curves, derivations, or coordinate vectors are naturally isomorphic.

\vspace{0.5cm}
\noindent
{\bf Part I}: Curves $\Longleftrightarrow$ coordinate vectors.

Take a chart around $p$:
$$
(x^1,\dots,x^n).
$$
Let $\gamma(t)$ be a curve with
$$
\gamma(0)=p.
$$
Define
$$
v^i
=
\frac d{dt}(x^i\circ \gamma)(0).
$$
Collect:
$$
v=(v^1,\dots,v^n)\in \mathbb R^n.
$$
So every curve determines a coordinate vector.

\vspace{0.5cm}
\noindent
{\em Independence of chart:}

Suppose another chart
$$
y^\alpha.
$$
Then by chain rule,
$$
\frac d{dt}(y^\alpha\circ\gamma)
=
\frac{\partial y^\alpha}{\partial x^i}
\frac d{dt}(x^i\circ\gamma).
$$
So components transform by Jacobian:
$$
v^\alpha
=
\frac{\partial y^\alpha}{\partial x^i}v^i.
$$
Exactly the transformation law for tangent vectors.
So this is geometric.

\vspace{0.5cm}
\noindent
{\em Converse:}

Given any
$$
(v^1,\dots,v^n),
$$
construct curve
$$
\gamma(t)
=
(x^1(p)+tv^1,\dots,x^n(p)+tv^n)
$$
in coordinates.
Its velocity is $v$.
So every coordinate vector comes from a curve.
Thus:
$$
\boxed{
\text{curves } \cong \mathbb R^n
}
$$

\vspace{0.5cm}
\noindent
{\bf Part II:} Curves $\Longrightarrow$ derivations.

Given curve $\gamma$, define:
$$
D_\gamma(f)
=
\frac d{dt}f(\gamma(t))\Big|_{0}.
$$
Linearity is obvious.
Leibniz rule:
$$
D_\gamma(fg)
=
\frac d{dt}(fg)(\gamma(t)).
$$
Product rule:
$$
f(p)D_\gamma(g)
+
g(p)D_\gamma(f).
$$
So it's a derivation.

\vspace{0.5cm}
\noindent
{\bf Part III:} Derivations $\Longrightarrow$ coordinate vectors.

Suppose
$$
D
$$
is a derivation.
Define components
$$
v^i:=D(x^i).
$$
Claim:
$$
D=
v^i
\frac{\partial}{\partial x^i}.
$$
Need to be proven.
For this, take smooth $f$ and Taylor expand near $p$:
$$
f
=
f(p)
+
(x^i-x^i(p))
\frac{\partial f}{\partial x^i}(p)
+
\text{quadratic}.
$$
Derivation kills constants:
$$
D(f(p))=0.
$$
Quadratic terms lie in $\mathfrak m^2$, derivations kill those.
So:
$$
D(f)
=
D(x^i)
\frac{\partial f}{\partial x^i}(p).
$$
Thus
$$
D
=
v^i \partial_i.
$$
Every derivation has coordinate components.

\vspace{0.5cm}
\noindent
{\em Converse:}

Given
$$
v^i
$$
define
$$
D
=
v^i\partial_i.
$$
Clearly derivation.
So
$$
\boxed{
\text{derivations}
\cong
\mathbb R^n
}
$$
Combining:
$$
\boxed{
\text{curves}
\cong
\text{derivations}
\cong
\mathbb R^n
}
$$
proved.

\vspace{0.5cm}
\noindent
{\em Embedded definition follows:}

For embedding
$$
r(u^1,\dots,u^m),
$$
pushforward gives
$$
dr_p\left(
\frac{\partial}{\partial u^i}
\right)
=
\frac{\partial r}{\partial u^i}.
$$
So embedded tangent vectors are just images of coordinate derivations.
Equivalent.

\vspace{0.5cm}
\noindent
{\em Implicit definition follows:}

If
$$
M={F=0},
$$
and $\gamma\subset M$,
then
$$
F(\gamma(t))=0.
$$
Differentiate:
$$
dF_p(\gamma'(0))=0.
$$
So
$$
T_pM\subseteq \ker dF_p.
$$
Dimensions agree, so
$$
T_pM=\ker dF_p.
$$
Equivalent.

\vspace{0.5cm}
\noindent
{\bf Elegant summary:}

All these are manifestations of first-order information at a point.
A tangent vector is:
\begin{itemize}
\item a first derivative of a curve,
\item a first-order differential operator,
\item a first-order coordinate increment.
\end{itemize}
All are the same notion.

This proof essentially shows:
$$
T_pM
\cong
(\mathfrak m_p/\mathfrak m_p^2)^*
$$
because derivations depend only on first-order Taylor terms.
That is why all this works.
It’s really a theorem about first-order infinitesimals.

%If you'd like, we can even prove why derivations annihilate (\mathfrak m^2), which is a beautiful little argument.


\end{document}
%
% ****** End of file apstemplate.tex ******


